\documentclass[11pt]{article}

\usepackage[a4paper,margin=0.60in]{geometry}
\usepackage{hyperref}
\usepackage{enumitem}
\usepackage{titlesec}
\usepackage{xcolor}
\usepackage{tabularx}
\usepackage{qrcode}

% ---- Colors ----
\definecolor{violettitle}{RGB}{124,58,237}
\definecolor{textgray}{RGB}{100,100,100}

\hypersetup{
	colorlinks=true,
	urlcolor=violettitle
}

% ---- Section titles ----
\titleformat{\section}
{\large\bfseries\color{violettitle}}
{}{0em}{}[\titlerule]

\setlength{\parindent}{0pt}

\setlist[itemize]{
noitemsep,
topsep=2pt,
itemsep=2pt,
leftmargin=1.4em
}

\begin{document}

%------------------------------------------------
% Header
%------------------------------------------------

\begin{center}

{\LARGE \textbf{Rimsha Afzal}} \\[6pt]

AI Developer — NLP, Machine Learning \& Generative AI \\[6pt]

Portfolio: \href{https://rimsha-afzal.github.io}{rimsha-afzal.github.io} \quad
GitHub: \href{https://github.com/rimsha-afzal}{github.com/rimsha-afzal} \quad
LinkedIn: \href{https://www.linkedin.com/in/rimsha-afzal-b0537123b/?lipi=urn%3Ali%3Apage%3Ad_flagship3_profile_view_base_contact_details%3BTbhXtoP0RS6goeD96Ggzjg%3D%3D}{rimsha-afzal} \\[4pt]
{\color{textgray}rimshaafzal2016@gmail.com $|$ +39 3913286575 $|$ Trento, Italy}

\end{center}

\vspace{-6pt}

%------------------------------------------------
% About Me
%------------------------------------------------

\section*{Profilo}

Studentessa magistrale in Artificial Intelligence Systems, con background in Computer Science ed esperienza pratica in NLP, sperimentazione di modelli di machine learning, transformer-based models e AI-driven decision-support pipelines. Ho lavorato su text preprocessing, preparazione di dataset, model fine-tuning, sentiment analysis, classificazione e valutazione dei modelli utilizzando Python, PyTorch e HuggingFace Transformers. Sono particolarmente interessata a ruoli di AI Developer focalizzati sulla costruzione, il testing e il miglioramento di sistemi AI applicati, con particolare attenzione a NLP, Generative AI, autonomous agents e soluzioni scalabili basate su modelli.

%------------------------------------------------
% Education
%------------------------------------------------

\section*{Formazione}

\textbf{Università di Trento, Italia} \hfill {\color{textgray}2025 -- In corso} \\[-1pt]
Laurea magistrale in Artificial Intelligence Systems \\[4pt]
{\color{textgray}Corsi rilevanti: Machine Learning, Natural Language Processing, Deep Learning, Autonomous Software Agents, Designing Large Scale AI Systems}

\vspace{8pt}

\textbf{Università di Trento, Italia} \hfill {\color{textgray}2021 -- 2025} \\[-1pt]
Laurea triennale in Computer Science \\[4pt]
{\color{textgray}Corsi rilevanti: Algorithms and Data Structures, Software Engineering, Human-Computer Interaction}

%------------------------------------------------
% Technical Skills
%------------------------------------------------

\section*{Competenze Tecniche}

\textbf{Programming \& Development} \\
{\color{textgray}Python, Java, Kotlin, HTML, CSS, Basic Scripting, Git}

\vspace{4pt}

\textbf{AI \& Machine Learning} \\
{\color{textgray}Machine Learning, Deep Learning, Natural Language Processing, Transformer Architectures, BERT, FinBERT, GPT-2, LoRA Fine-Tuning, Language Modeling}

\vspace{4pt}

\textbf{Frameworks \& Tools} \\
{\color{textgray}PyTorch, HuggingFace Transformers, HPC/Slurm, Figma}

\vspace{4pt}

\textbf{Data \& Experimentation} \\
{\color{textgray}Data Preprocessing, Text Preprocessing, Tokenization, Train/Test Splitting, Model Evaluation, Metric Definition, Decision-Support Pipelines}

\vspace{4pt}

\textbf{Product \& Collaboration} \\
{\color{textgray}Requirement Translation, Workflow Analysis, Technical Documentation, Stakeholder Interviews, MVP Definition, Feature Testing}

%------------------------------------------------
% Projects
%------------------------------------------------

\section*{Progetti AI \& Tecnici}

\textbf{NLP-Based Financial Complaint Prioritization System} \hfill {\color{textgray}Tesi triennale} \\[-0.8em]
\begin{itemize}
\item Progettato e valutato una \textbf{NLP-based decision-support pipeline} per classificare reclami finanziari e supportare la prioritizzazione dei casi più urgenti
\item Elaborato un dataset reale di reclami attraverso preparazione dei dati, text preprocessing, suddivisione train/test e tokenization
\item Effettuato il fine-tuning di modelli \textbf{BERT} e \textbf{FinBERT}, utilizzando rappresentazioni transformer-based per attività di classificazione e sentiment analysis
\item Confrontato gli output dei modelli per identificare segnali utili all’analisi dei reclami e alla stima del livello di urgenza
\item Combinato risultati di classificazione e sentiment analysis in una \textbf{urgency metric}, collegando le predizioni dei modelli a processi decisionali operativi
\end{itemize}

\vspace{4pt}

\textbf{Transformer Language Model Optimization} \hfill {\color{textgray}Progetto in corso} \\[-0.8em]
\begin{itemize}
\item Sperimentazione con \textbf{GPT-2 tuning} e \textbf{LoRA-based fine-tuning} per comprendere come adattare language models a casi d’uso pratici
\item Esplorazione di tecniche di model adaptation, evaluation e optimization per transformer-based language models
\item Sviluppo di competenze pratiche su Generative AI workflows, comportamento dei modelli e strategie di fine-tuning
\end{itemize}

\vspace{4pt}

\textbf{Planning-Based Autonomous Software Agents} \hfill {\color{textgray}Progetto in corso} \\[-0.8em]
\begin{itemize}
\item Studio di agentic systems per attività di planning, belief management e decision-making in ambienti dinamici
\item Analisi di come autonomous software agents possano ragionare su obiettivi, gestire stati interni e supportare l’esecuzione di task
\item Applicazione di concetti di AI planning e intelligent systems per comprendere la progettazione di goal-driven AI applications
\end{itemize}

\vspace{4pt}

\textbf{OpenMove Internship — Internal Support Solution Prototype} \hfill {\color{textgray}Internship} \\[-0.8em]
\begin{itemize}
\item Condotto stakeholder interviews per comprendere workflow di customer care, criticità operative e bisogni degli utenti
\item Tradotto insight qualitativi in functional requirements, flussi di navigazione, logiche delle schermate e decisioni di prodotto per una nuova internal support solution
\item Progettato Figma mockups e un interactive prototype per validare chiarezza dei workflow, usabilità e copertura dei requisiti
\item Sviluppato esperienza nella traduzione di bisogni utente e business in specifiche tecniche strutturate, bilanciando usabilità e fattibilità
\end{itemize}

\vspace{4pt}

\textbf{Start-up Lab — MVP Definition \& Technical Feasibility} \hfill {\color{textgray}Team Project} \\[-0.8em]
\begin{itemize}
\item Guidato un team project dall’esplorazione del problema e dalla validazione dell’idea fino alla definizione dell’MVP
\item Condotto interviste esterne per testare ipotesi, raccogliere feedback e valutare la rilevanza del problema identificato
\item Tradotto gli insight di ricerca in un concept di soluzione concreto, contribuendo alla definizione di scope, user value e prime priorità tecniche
\item Preparato presentazioni per spiegare product logic, percorso di sviluppo, potenziale di business e considerazioni di fattibilità
\end{itemize}

%------------------------------------------------
% Community Engagement
%------------------------------------------------

\section*{Community Engagement}

\textbf{Italian Tech Week Volunteer} \hfill {\color{textgray}Oct 2025} \\
{\color{textgray}Supporto alla gestione operativa dell’evento e alla registrazione dei partecipanti durante una conferenza internazionale di tecnologia della durata di tre giorni.}

%------------------------------------------------
% Languages
%------------------------------------------------
\section*{Lingue}

\vspace{-2pt}

\begin{tabularx}{\textwidth}{@{}Xr@{}}
\begin{minipage}[t]{0.55\textwidth}
Inglese — {\color{textgray}C1} \\[2pt]
Italiano — {\color{textgray}C1} \\[2pt]
Urdu — {\color{textgray}B2}
\end{minipage}
&
\begin{minipage}[t]{0.25\textwidth}
\raggedleft
\vspace{-8pt}
{\color{textgray}\textbf{Portfolio}} \\[3pt]
\qrcode[height=1.9cm]{https://rimsha-afzal.github.io}
\end{minipage}
\end{tabularx}

\end{document}